% Copy-ready Section VIII insert for the post-closure Version 2 manuscript.
% Do not include this file from main.tex until the attained-demand reindexing
% and pointwise aggregate weight identity have been checked.

\subsection{The all-deficit partition function}
\label{subsec:all-deficit-partition-function}

Fix a block-level matching support $P$.  For every selected cell $e\in P$, let
$s_e,t_e$ be its endpoint block sizes and put
\[
  m_e:=\min\{s_e,t_e\},
  \qquad
  d_e:=|s_e-t_e|.
\]
If the actual high multiplicity is $j_e$, write
\[
  j_e=m_e-h_e.
\]
The canonical high condition implies $2h_e<m_e$.

After summing the literal partial-stub-matching fibre in every selected cell,
the aggregate bare weight is
\begin{equation}
  w(P,j)
  =
  \frac{\displaystyle\prod_{e\in P}
    (s_e)_{j_e}(t_e)_{j_e}}
  {\displaystyle (n)_J\prod_{e\in P}j_e!}
  \prod_{e\in P}g(j_e),
  \qquad
  J:=\sum_{e\in P}j_e.
  \label{eq:aggregate-bare-weight-v2}
\end{equation}
No objectwise full completion is selected in
\eqref{eq:aggregate-bare-weight-v2}; the formula is the sum of the complete
finite physical fibre.

For endpoint sizes $m,m+d$, define
\begin{equation}
  R_{m,d}(h)
  :=
  \frac{\binom mh}{(d+1)(d+2)\cdots(d+h)}
  2^{-hm+h(h+1)/2},
  \label{eq:local-deficit-ratio-v2}
\end{equation}
with $R_{m,d}(0)=1$.

\begin{lemma}[Aggregate deficit comparison]
\label{lem:aggregate-deficit-comparison-v2}
For a fixed block support $P$ and an admissible deficit vector
$h=(h_e)_{e\in P}$,
\[
  \frac{w(P,m-h)}{w_{\mathrm{full}}(P)}
  \le
  \prod_{e\in P}n^{h_e}R_{m_e,d_e}(h_e).
\]
\end{lemma}

\begin{proof}
The ratio of the local factors in one selected cell is exactly
\eqref{eq:local-deficit-ratio-v2}.  Put
\[
  H:=\sum_{e\in P}h_e.
\]
The only nonlocal change is the ambient falling-factorial denominator:
\[
  \frac{(n)_{J+H}}{(n)_J}
  =(n-J)_H
  \le n^H.
\]
Distributing $n^H=\prod_e n^{h_e}$ over the selected cells gives the result.
\end{proof}

\begin{lemma}[Cellwise optional-deficit product]
\label{lem:cellwise-deficit-product-v2}
Let $A_e$ be the finite positive-deficit set in cell $e$, and let
$u_e(h)\ge0$ be its charged local weight.  Then
\[
  \sum_{\text{optional deficit choices }\omega}
    \prod_{e\in P}u_e(\omega_e)
  =
  \prod_{e\in P}
    \left(1+\sum_{h\in A_e}u_e(h)\right),
\]
where the local factor is one if no positive deficit is selected.
Consequently, if
\[
  \sum_{h\in A_e}u_e(h)\le \sigma_e
  \qquad(e\in P),
\]
then the total optional-deficit weight is at most
\[
  \prod_{e\in P}(1+\sigma_e).
\]
\end{lemma}

\begin{proof}
Expand the finite product.  Each monomial selects either the constant term one
or exactly one positive-deficit weight in every distinguishable selected cell.
\end{proof}

For $2h<m$, the integer estimate
\[
  h\left\lfloor\frac{2m}{3}\right\rfloor
  \le
  hm-\frac{h(h+1)}2
\]
and \eqref{eq:local-deficit-ratio-v2} give
\begin{equation}
  n^hR_{m,d}(h)
  \le
  \left(
    \frac{nm}{2^{\lfloor2m/3\rfloor}}
  \right)^h.
  \label{eq:two-thirds-local-charge-v2}
\end{equation}
Define the cell-dependent charge
\[
  \rho_e
  :=
  \frac{n m_e}{2^{\lfloor2m_e/3\rfloor}}.
\]
For all sufficiently large $n$, every endpoint charge is at most $1/2$.
Therefore
\[
  \sum_{h\in A_e}n^hR_{m_e,d_e}(h)
  \le
  \sum_{h\ge1}\rho_e^h
  =
  \frac{\rho_e}{1-\rho_e}
  \le2\rho_e.
\]
By Lemma~\ref{lem:cellwise-deficit-product-v2},
\begin{equation}
  \sum_{h:\,m-h\text{ high}}w(P,m-h)
  \le
  w_{\mathrm{full}}(P)
  \prod_{e\in P}(1+2\rho_e).
  \label{eq:fixed-support-all-deficit-v2}
\end{equation}

Let
\[
  \rho_n:=\max_e\rho_e.
\]
The four-size phase relation gives
\[
  \rho_n
  =
  O\!\left(
    \frac{(\log n)^{7/3}}{n^{1/3}}
  \right)
  =o(1).
\]
Since $|P|\le k_{\mathrm{co}}=\Theta(n/\log n)$,
\[
  \prod_{e\in P}(1+2\rho_e)
  \le
  (1+2\rho_n)^{k_{\mathrm{co}}}
  \le
  \exp\!\left\{
    O\!\left(n^{2/3}(\log n)^{4/3}\right)
  \right\}.
\]
This estimate has no extra factor of order $\log n$ from counting the
admissible deficits.

Combining \eqref{eq:fixed-support-all-deficit-v2} with the exact full-endpoint
normalisation, endpoint transportation, and the partial-diagonal estimate gives
\begin{equation}
  \operatorname{BareSkeletonSum}_n
  \le
  \exp\!\left\{
    O\!\left(n^{2/3}(\log n)^{4/3}\right)
    +O(\sqrt{n\log n})
  \right\}
  =
  \exp\!\left\{
    o\!\left(\frac{n}{(\log n)^4}\right)
  \right\}.
  \label{eq:bare-skeleton-sharp-v2}
\end{equation}

% Audit boundary: equation \eqref{eq:bare-skeleton-sharp-v2} may be promoted
% into the canonical manuscript only after the attained-demand reindexing and
% the pointwise aggregate weight identity are proved on the integrated branch.
